\section{Bargaining over development opportunities}
\label{sec:bargaining}

\subsection{Responses differ with the exchange and its recipient}

The 526 negotiation threads contain 920 responses: 113 acceptances, 394 counters and 413 rejections. Acceptance is legal in 918 opportunities; two would leave a party cash-negative after mandatory charges. Table~\ref{tab:offers} describes whose color group the offered exchange would complete.

\begin{table}[t]
\centering\small
\begin{tabular}{@{}lrlrlrl@{}}
\toprule
& \multicolumn{2}{c}{All offers} & \multicolumn{2}{c}{Opening offers} & \multicolumn{2}{c}{Later offers}\\
\cmidrule(lr){2-3}\cmidrule(lr){4-5}\cmidrule(lr){6-7}
Newly completes & $N$ & A/C/R & $N$ & A/C/R & $N$ & A/C/R\\
\midrule
Author only & 289 & 19/111/159 & 187 & 3/46/138 & 102 & 16/65/21\\
Recipient only & 140 & 24/78/38 & 51 & 2/18/31 & 89 & 22/60/7\\
Both & 81 & 10/35/36 & 48 & 3/11/34 & 33 & 7/24/2\\
Neither & 408 & 60/169/179 & 238 & 24/62/152 & 170 & 36/107/27\\
\midrule
All & 918 & 113/393/412 & 524 & 32/137/355 & 394 & 81/256/57\\
\bottomrule
\end{tabular}
\caption{All acceptance-legal offers, separated by position in the thread. A/C/R are counts of acceptances, counters and rejections. Completion concerns the current author's and recipient's holdings, including mortgaged groups. Splitting by position exposes composition differences; assets, prices, recipients and game states still differ within each cell.}
\label{tab:offers}
\end{table}

Author-only completion offers are rejected in 159/289 cases (55.0\%), compared with 38/140 (27.1\%) for recipient-only completion. Much of this pooled contrast changes when negotiation stage is separated. Among opening offers, rejection counts are 138/187 (73.8\%) versus 31/51 (60.8\%); among later offers, they are 21/102 (20.6\%) versus 7/89 (7.9\%). Author-only offers more often open a thread, where rejection is common in both categories. The within-position differences are about thirteen percentage points, compared with twenty-eight points pooled. These are conditional descriptions, not an identified effect of completion.

The response contrast also varies by recipient. Pooled over thread position, Sol, Gemini and Grok reject author-only offers more often than recipient-only offers and accept them less often. Opus shows the opposite pooled orderings, rejecting 3/52 versus 6/42 and accepting 7/52 versus 5/42, while mostly countering both. Within opening and later offers separately, the rejection ordering holds in five of the six Sol/Gemini/Grok cells, while the acceptance ordering holds in three, ties in one and reverses in two. The appendix preserves these breakdowns. The descriptive pattern is therefore most consistent for rejection, not a uniform preference across every configuration and stage.

\begin{table}[t]
\centering\small
\begin{tabularx}{\linewidth}{@{}lrrrrY@{}}
\toprule
Recipient & Received & Accept & Counter & Reject & Response condition\\
\midrule
Sol & 235 & 30 & 140 & 65 & Standard\\
Gemini & 242 & 26 & 36 & 180 & Standard\\
Opus & 189 & 29 & 141 & 19 & Action-only in 162 responses\\
Grok & 254 & 28 & 77 & 149 & Eleven text/call conflicts\\
\bottomrule
\end{tabularx}
\caption{All 920 responses, including two without a legal acceptance. No response is an engine default. Seven of Grok's acceptances, one counter and three rejections conflict with the action named in its content. Counts describe received responses under these conditions.}
\label{tab:styles}
\end{table}

Each configuration accepts twenty-six to thirty offers, but disagreement takes different forms (Table~\ref{tab:styles}). Sol and Opus mostly counter; Gemini and Grok mostly reject. Only thirty-two settlements occur on an opening offer, with eighty-one following a counter. This is compatible with continued bargaining in selected negotiable threads, not evidence that introducing a counter would improve an otherwise identical negotiation.

Initiation differs more. Sol opens 132 threads, Gemini eighteen, Opus 240 and Grok 136. The response condition changes the interpretation of the last count: sixty-four Grok threads begin with a received proposal whose accompanying content names another action, and seventeen of those threads settle. Seven of Grok's twenty-eight acceptances similarly conflict with content; three complete Opus groups. They remain actual proposals and settlements that other participants encounter. They are not clean measures of Grok's intention to negotiate or agree.

Low acceptance does not rule out concessions on poor terms; judging agreement bias would require a utility or fair-price reference.

\FloatBarrier

\subsection{Counters reprice exchanges with and without completion}

Counter contents show how disagreement changes a transaction. Of 111 counters to author-only completion offers, ninety-two retain every requested deed and preserve the original author's completion. Those two predicates select the same records here. Eighty-eight still leave only that party completing. Of those eighty-eight, eighty-five improve the counterer's net cash; the median cash change across the eighty-eight is an increase of \$125. Seventy-nine preserve the entire noncash bundle, including deeds and jail cards, and all seventy-nine increase the counterer's cash. These counters negotiate a price for the same exchange rather than removing the completing deed.

This preservation and repricing is also common without completion (Table~\ref{tab:counterbaseline}). Among 110 counters to non-completing offers whose authors request at least one deed, eighty-seven retain every requested deed and eighty-one preserve the entire noncash bundle while increasing the counterer's cash. Sol and Opus write ninety-two of the 111 counters and eighty of the 110, consistent with their greater use of counteroffers. Repricing the unchanged bundle occurs in a majority of each configuration's counters in both populations. It is an observed bargaining practice in these games, not evidence of a special completion premium.

\begin{table}[H]
\centering\small
\begin{tabularx}{\linewidth}{@{}Yrrr@{}}
\toprule
Initial offer completes & Counters & \shortstack{All requested\\deeds retained} & \shortstack{Same noncash bundle,\\more cash to counterer}\\
\midrule
Author only & 111 & 92 & 79\\
Neither party & 110 & 87 & 81\\
\bottomrule
\end{tabularx}
\caption{Counters to offers requesting at least one deed. The same deed-retention predicate is used in both rows. Noncash equality includes deeds and jail cards on both sides. The second row is a descriptive comparison drawn from 169 non-completing offers that received counters; fifty-nine requested no deed. It is not a matched control.}
\label{tab:counterbaseline}
\end{table}

An exploratory comparison with board prices gives no consistent higher ratio for bids completing the bidder's group across the position and mortgage-status cuts. Seller asks are higher for recipient-completing deeds in those cuts, but neither deed composition nor state is held fixed (Appendix~\ref{app:prices}). Board price is not a valuation standard.

Explicit explanations supply evidence about what participants consider. Both LLM-assisted readings identify a stated trade-linked rival consequence in sixty-one of seventy-four outright rejections that would complete the recipient's group. The agreed labels are Gemini 29/31, Grok 27/31, Sol 4/4 and Opus 1/8; all eight Opus records lack a private note. They occur in 34/36 both-completing offers and 27/38 recipient-only offers. The reasons concern rival ownership, building funds and survival, including cases where the rival receives no new group. Thirty-four records come from G1 and G15, so this is concentrated textual evidence rather than a balanced comparison of configurations. The following exchanges show the strategic considerations behind the labels.

Of the nineteen accepted author-only completion offers, sixteen follow a counter, eight recipients hold less than \$100, and seven recipients are Opus under action-only recovery. These categories overlap. Three acceptances are Grok text/call conflicts. All three opening-offer acceptances are Grok's: G16 d247 and G18 d083 conflict, while G3 d194 names acceptance in both content and call.

\TranscriptTitle{g2}{Pricing a rival's completion, then buying the group instead}

At G2 turn 45, Opus has two yellows and offers Sol \$325 for Marvin Gardens. Sol asks \$650:
\TranscriptQuote{g2-d124-price}
\TranscriptQuote{g2-d124-note}
Sol's ask also changes with the buyer's position. One turn earlier, it countered Grok's \$200 bid for Marvin Gardens at \$450; Grok then owned no yellow. Against Opus, which would complete yellow, Sol asks \$650, reduces this to \$600 and refuses \$450. This is a same-seller, same-deed contrast in stated prices, although counterparty and negotiation order also change. Sol's note expressly values the rival's development capacity. Its \$520 residual-cash calculation is correct, but that cash buys three yellow houses, not the nine claimed in the quote.

Opus calls a final price while keeping a higher private ceiling:
\TranscriptQuote{g2-d127-final}
\TranscriptQuote{g2-d127-ceiling}
Sol refuses:
\TranscriptQuote{g2-d128-keep}
At turn 69, the direction reverses: Sol has \$435 and seeks Opus's two yellows in return for two railroads. Opus now has \$17 and replies:
\TranscriptQuote{g2-d197-reverse}
\TranscriptFinding{Sol accepts the added \$100 and completes yellow. Opus is eliminated three turns later, owing Gemini \$1,300, and its newly acquired railroads pass to Gemini. Sol spends \$300 on yellow houses, later receives \$150 salvage, and collects no developed yellow rent. It does collect \$44 of undeveloped Atlantic rent at turn 77. Its later bankruptcy also transfers the yellows to Gemini, the eventual winner.}
\TranscriptReading{The parties bargain over the rival's development budget, a stated price limit and the value of the completing deed. Their roles then reverse as Opus's cash falls. The group ultimately reaches a third player through bankruptcy, linking the bilateral negotiations to the wider competitive economy.}
\TranscriptSource{G2 d119--128 and d197--198. Public quotations appear in subsequent counterparty requests. The \$44 receipt is event 1566; developed rent and salvage are separate measures.}

\TranscriptTitle{g15}{Refusing one's own complete group three times}

At G15 turn 63, Sol has \$730 and Grok \$333. Sol offers its two greens for Grok's Ventnor and \$100. Both would complete a group:
\TranscriptQuote{fx-g15-offer}
Grok's refusal focuses on the unequal cash left for construction:
\TranscriptQuote{fx-g15-refusal-note}
The proposed balances, \$830 for Sol and \$233 for Grok, agree with the ledger. One house on each green costs \$600, as the note says. That is not the minimum legal construction: one green house costs \$200. The explanation concerns a full first layer and leaves other financing possibilities untested.

Sol removes its cash request and adds Kentucky, then makes a third offer:
\TranscriptQuote{fx-g15-sweetener}
\TranscriptQuote{fx-g15-final-refusal}
\TranscriptFinding{No exchange settles. Eighteen turns later Grok is the first player eliminated, and its yellow and green blockers pass to Gemini, which wins. Sol never completes either group.}
\TranscriptReading{Both parties discuss relative ability to develop. Sol changes the offer to address Grok's stated cash objection, but Grok still refuses its own completion. The disagreement concerns the position each player could finance after the swap, rather than ownership completion alone.}
\TranscriptSource{G15 d168--175, standard first attempts. The quoted public messages are present in the counterparties' next actual requests. Cash consequences come from the offer ledger.}

\subsection{Cash at settlement and bargaining limits}

An exploratory cut describes the holders who release completing deeds. Among forty-eight cash settlements completing the paying party's group, median seller cash is \$194, compared with \$571 in fifty-seven cash settlements completing neither party's group. Twenty-one of those forty-eight sellers hold less than \$100. Conversely, among thirty cash-receiving settlements with a seller below \$100, twenty-one complete the payer's group. The latter denominator conditions on low cash rather than on completion.

The association appears across the reported \$50, \$100, \$150 and \$200 thresholds (Appendix~\ref{app:prices}). Opus supplies thirteen of the twenty-one low-cash sellers of completing deeds, and six of the twenty-one eventually win. These settlements illustrate a recurring tradeoff: releasing a rival's completing deed supplies immediately usable cash. Their selected prices and circumstances do not identify the effect of cash scarcity on agreement.

\TranscriptTitle{g7}{A settled public bluff about the final price}

Sol asks \$625 for North Carolina. Grok counters:
\TranscriptQuote{g7_d42_public}
Its private note on the same response retains a higher contingency:
\TranscriptQuote{g7_d42_private}
Sol answers:
\TranscriptQuote{g7_d43_public}
\TranscriptFinding{The deed transfers for \$600. Sol's note prefers reinvesting to holding out for another \$25. Grok completes green through a \$171 Pennsylvania auction purchase at turn 95, but never builds on the group. At turn 110 it is eliminated owing Sol \$1,400. North Carolina returns to Sol within the complete green group, and Sol builds a house on each deed two turns later.}
\TranscriptReading{The public final price is below the privately recorded contingency: an authorized bargaining bluff followed by agreement, contrasting with G2's refused final offer. The later estate returns the same asset to its original seller as part of a complete group. The record shows both the tactic and the succession of owners, without isolating the bluff's effect on the price.}
\TranscriptSource{G7 d41--43, standard first attempts; acquisition event 1680, estate events 1946--1954 and build events 1971--1973. Sol receives the \$600 message.}

\TranscriptTitle{g19}{A false premise corrected while the deal is accepted}

Grok's received call proposes Virginia plus \$150 for States:
\TranscriptQuote{g19_d135_false_public}
Both deeds are pink. The same response's private note says:
\TranscriptQuote{g19-abort}
Assistant content names \End, but the received tool proposes the exchange. Sol responds:
\TranscriptQuote{g19_d136_correction_accept}
\TranscriptFinding{The swap and \$150 payment execute and complete no group. Sol's note calls the exchange a ``strictly favorable swap.''}
\TranscriptReading{Sol separates the offered terms from the false explanation: it corrects the color claim and accepts the cash-plus-deed exchange. Grok's public message supports the proposal, while its content and attached note oppose it. This is a useful contrast to the bluffs: accepting an offer need not mean accepting the proposer's account of why it is attractive.}
\TranscriptSource{G19 d135--136, standard first attempts. The false claim is present in Sol's delivered request.}
